\providecommand{\toplevelprefix}{../..}  %
\documentclass[../../book-main.tex]{subfiles}

\begin{document}

\chapter*{Preface to Version 2.0}

% \begin{center}
%     ``{\em All roads lead to Rome}.''

% \end{center}
% \vspace{5mm}

As an open-sourced book, the first version has been frequently and constantly updated in the past six months since its release on August 18, 2025. Nevertheless, based on our experience from using the book for a new course taught at the University of Hong Kong in Fall 2025 and feedback we have received from our colleagues, teaching assistants and students, we had noticed several issues with the previous version that cannot be easily addressed with incremental changes. 

Hence, we have decided to make some substantial changes and upgrades of the  content and organization of the book with a new version. The new version also allows us to explicitly reveal strong conceptual and technical connections among materials across different chapters and sections, so as to improve the overall pedagogical values of the book. 

In our opinion, the new version gives a much more unified and stream-lined presentation of this subject that it deserves. The new version also reflects the quickly improving theoretical understanding of this subject as well as an explosive number of new practical applications of learning deep representations for real-world data, that have not been systematically documented anywhere else.

\subsection*{Major Changes in Version 2.0}
\begin{itemize}
    \item We have split Chapter 3 in the first version into to two chapters, now Chapter \ref{ch:denoising} and \ref{ch:lossy-compression}. The new Chapter \ref{ch:denoising} focuses on the denoising process for learning low-dimensional distributions. We add a new section that characterizes conditions under which the process leads generalization or memorization. The new Chapter \ref{ch:lossy-compression} focuses on representation learning based on a lossy coding approach and the principle of maximizing information gain. Since the old version only illustrated how to apply the principle to learn image representations in the supervised setting, the new version has added a case study with the unsupervised setting. It also provides theoretical justification for popular unsupervised learning methods such as constrative learning and DINO in particular.
    \item In the new Chapter \ref{ch:deep-networks} on designing deep representations via unrolled optimization, we have added a principled  derivation of a causal version for the white-box transformer architecture CRATE, which is important for processing sequential data such as texts as we often see in applications.
    \item In the new Chapter \ref{ch:consistent-representations} on consistent representation learning, we added a new section that further elucidate the relationships between distribution learning and representation learning. This provides theoretical clarification and justification for popular practical autoencoding methods such as VAE and RAE. 
    \item In the new Chapter \ref{ch:conditional-inference}, we have added a section that provides principled explanation to representation learning with paired data and  conditioned generation through the perspective of mutual information. It provides theoretical justification for popular practical methods such as CLIP.
    \item The new application Chapter \ref{ch:applications} now features many more detailed implementations of representation learning for real-world data, including natural 2D images,  natural 3D objects, human body motions, and natural languages, under almost all popular and practical settings, supervised, weakly supervised, unsupervised, and conditioned. 
    \item The final Chapter \ref{ch:future} about open directions has also been significantly rewritten. We attempt to give a clearer taxonomy of different levels of intelligence so that we can clarify what we have done and understood about intelligence (with this book) and what have not. We believe that open problems associated with more advanced forms of intelligence can be so well posed that they can be studied qualitatively and quantitatively via scientific and mathematical means, instead of remaining as a mysterious subject.
\end{itemize}


\subsection*{Content Contributors for Version 2.0.}
Besides the authors, many students and colleagues have started to join this project and contributed valuable content to different parts of the book, during the preparation of {\em Version 2.0} of the book. Below is an incomplete list of people and their specific contributions, in alphabetical order:
\begin{itemize}
    \item Tianzhe Chu: Experiments in Section \ref{sec:image_classification}, and AI assistant for the book.
    \item Prof. Shenghua Gao: Section \ref{sec:Michelangelo} and \ref{sec:CUPID}.
    \item Bingbing Huang: Sections \ref{sec:Michelangelo} and \ref{sec:CUPID}.
    \item Prof. Qing Qu: Section \ref{sec:diffusion_memorization}.
    \item Shengbang Peter Tong: Sections \ref{sec:continuous} and  \ref{sec:RAE-implement}.
    \item Chengyu Wang: Sections \ref{sec:Michelangelo} and \ref{sec:CUPID}.
    \item Ziyang Robin Wu: Section \ref{sec:CLIP}.
    \item Chun-Hsiao Daniel Yeh: Sections \ref{sec:auto-encode-img-gen} and \ref{sec:cond_gen}.
    \item Brent Yi: Section \ref{sec:egoallo}.
    \item Jingfeng Yang: Section \ref{sec:CLIP}.
    \item Dr. Yaodong Yu: Chapter \ref{ch:deep-networks} is based on his PhD thesis.
    \item Dr. Zibo Zhao: Section \ref{sec:Michelangelo}.

\end{itemize}

\end{document}